
\documentclass[aspectratio=169]{beamer}

\mode<presentation>
\usetheme{Boadilla}
\definecolor{NTHU1}{RGB}{127,16,132}
\definecolor{NTHU2}{RGB}{228,210,231}
\definecolor{blue}{RGB}{30,90,205}
\definecolor{red}{RGB}{213,94,0}
\definecolor{green}{RGB}{0,128,0}
\definecolor{charcoalgray}{RGB}{108,104,104}
\setbeamercolor{title}{fg=NTHU1}
\setbeamercolor{frametitle}{fg=NTHU1}
\setbeamercolor{block title}{bg=NTHU1, fg=white}
\setbeamercolor{block body}{bg=NTHU2}
\setbeamercolor{structure}{fg=NTHU1}
\setbeamercolor{item projected}{fg=white}
\setbeamercolor{item}{fg=NTHU1}
\setbeamercolor{subitem}{fg=NTHU1}
\setbeamercolor{section in toc}{fg=NTHU1}
\setbeamercolor{description item}{fg=NTHU1}
\setbeamercolor{caption name}{fg=NTHU1}
\setbeamercolor{button}{bg=NTHU1, fg=white}
\setbeamercolor{caption name}{fg=NTHU1}
\usepackage{graphics}
\usepackage{tikz}
\usepackage{amsmath}
\usepackage{bbm}
\usetikzlibrary{decorations.pathreplacing}
\usepackage{geometry}
\usepackage{booktabs}
\usepackage{bookmark}
\usepackage{multirow, makecell}
\usepackage{float}
\usepackage{fancyvrb}
\usepackage{caption}
\captionsetup[figure]{singlelinecheck = false, format= hang, justification=centering}
\captionsetup[subfigure]{singlelinecheck = false, format= hang, justification=raggedright}
\usepackage{subcaption}
\usepackage{adjustbox}
\usepackage{hyperref}
\usepackage{threeparttable}
\usepackage{appendixnumberbeamer} 
\usepackage[T1]{fontenc}
\usepackage[default]{lato} 
\usepackage{smartdiagram}
\smartdiagramset{
  back arrow disabled = true,
  module minimum width=1cm,
  module x sep = 3,
  uniform arrow color=true,
}
\usepackage{xcolor}
\usepackage{colortbl}
	\definecolor{light-gray}{gray}{0.99}

\newenvironment{wideitemize}{\itemize\addtolength{\itemsep}{10pt}}{\enditemize}
\newenvironment{wideenumerate}{\enumerate\addtolength{\itemsep}{10pt}}{\endenumerate}
\newenvironment{widedescription}{\description\addtolength{\itemsep}{10pt}}{\enddescription}
\hypersetup{
colorlinks=true,
linkcolor=NTHU1,
filecolor=green, 
urlcolor=blue,
}
\beamertemplatenavigationsymbolsempty
\setbeamercolor{author in head/foot}{bg=white, fg=NTHU1}
\setbeamercolor{title in head/foot}{bg=white, fg=NTHU1}
\setbeamercolor{date in head/foot}{bg=white, fg=NTHU1}
\setbeamercolor{section in head/foot}{bg=white, fg=NTHU1}
\setbeamercolor{headline}{bg=white}
\setbeamertemplate{footline}{
    \leavevmode%
    \hbox{%
        \begin{beamercolorbox}[wd=.333333\paperwidth,ht=2.25ex,dp=1ex,center]{date in head/foot}%
            \usebeamerfont{date in head/foot}%\insertdate
        \end{beamercolorbox}%
        \begin{beamercolorbox}[wd=.333333\paperwidth,ht=2.25ex,dp=1ex,center]{title in head/foot}%
            \usebeamerfont{title in head/foot}\insertshorttitle
        \end{beamercolorbox}%
        \begin{beamercolorbox}[wd=.333333\paperwidth,ht=2.25ex,dp=1ex,right]{date in head/foot}%
            \usebeamerfont{date in head/foot}\insertframenumber{}\hspace*{2em}
        \end{beamercolorbox}}%
        \vskip0pt%
    }
    %% May need to script the introduction!!!!!!
%\setbeamercolor{page number in head/foot}{fg=NTHU1}
\setbeamertemplate{section in toc}[sections numbered]
\setbeamertemplate{subsection in toc}{\leavevmode\leftskip=3em\rlap{\hskip-1.75em\inserttocsectionnumber.\inserttocsubsectionnumber}
\inserttocsubsection\par}
\setbeamerfont{subsection in toc}{size=\footnotesize}
%\setbeamertemplate{headline}{%
  %\begin{beamercolorbox}[ht=5.5ex]{section in head/foot}
    %\vskip2pt\insertnavigation{0.33\paperwidth}\vskip2pt
  %\end{beamercolorbox}%
%}
\newenvironment{transitionframe}{\setbeamercolor{background canvas}{bg=white}\setbeamertemplate{footline}{} \begin{frame}[noframenumbering]}{\end{frame}}


\makeatletter
\let\@@magyar@captionfix\relax
\makeatother


\title[]{Public Economics and Development} % Change this regularly
\subtitle[]{Week 12: Allocation of Authority Across Bureaucrats and Managers, and Corruption}
\author[Lee]{Seung-hun Lee }
\institute[]{Taipei School of Economics\\ National Tsing Hua University}

\date[]{May 12th, 2026}

\begin{document}
\begin{transitionframe}
\titlepage
\end{transitionframe}


%%% Color slides for section headers: Use for colloquium version (The ones bewteen \iffals and \fi)

\section{Overview setup of the topic}

\begin{frame}
\frametitle{Why does the public sector personnel matter?} 
Take up huge roles: In principle and in data
\begin{itemize}\setlength\itemsep{3pt}
\item As in other sectors, they embody human capital of the state
\item They determine provision of public services and implementation of policies
\item Takes up nearly 20\% of all employment  and 35\% of formal employment (World Bank)
\item Responsible for the public sector that takes up approximately 45\% of GDP (OECD)
\end{itemize}\medskip
However, the quality of public sector workers varies across countries
\begin{itemize}\setlength\itemsep{3pt}
\item 20-35\% of front-line workers (teachers, health care) absent in developing countries
\item Lack institutional structures to incentivize, monitor existing workers
\begin{itemize}\setlength\itemsep{3pt}
\item Poor structure $\to$ Absenteeism and inefficiencies $\to$ Bad public services (repeat)
\end{itemize}
\item Finding motivated and capable individuals to work is difficult
\begin{itemize}\setlength\itemsep{3pt}
\item Lower pay, more rigid mobility than private sector
\end{itemize}
\item Bureaucratic quality is also clustered with development 
\end{itemize}
\end{frame}

\begin{frame}
\frametitle{Public sector takes up huge share of employment globally} 
\begin{figure}
  \includegraphics[keepaspectratio, width=0.8\textwidth]{fig1.png}
\caption{Share of public sector employment in total employment, 2019}
\end{figure}
\begin{footnotesize}
Source: Our World in Data (\hyperlink{https://ourworldindata.org/grapher/public-sector-employment-as-a-share-of-total-employment?time=2019}{link}), reprocessed from World Bank and ILO data 
\end{footnotesize}
\end{frame}

\begin{frame}
\frametitle{Clustering of development and bureaucratic quality} 
\begin{figure}
  \includegraphics[keepaspectratio, width=0.6\textwidth]{fig2.png}
\end{figure}
\begin{footnotesize}
Source: Besley, Burgess, Khan, Xu (2022) ``Bureaucracy and Development", \textit{Annual Review of Economics}, 14, 397-424 
\end{footnotesize}
\end{frame}




\begin{frame}
\frametitle{Corruption: A Global Development Challenge}
Corruption is rampant — especially in developing countries
\begin{itemize}\setlength\itemsep{3pt}
\item Over 2/3 of countries score below 50 on the Corruption Perceptions Index (CPI)
\begin{itemize}\setlength\itemsep{3pt}
\item CPI ranges from 0 (highly corrupt) to 100 (very clean)
\item Global average CPI score stuck at 43 for over a decade — little progress
\item 148 out of 180 countries have stagnated or worsened since 2012
\end{itemize}
\item Undermines public service delivery, economic growth, and trust in government
\item Diverts public resources away from those who need them most
\end{itemize}\medskip
Yet research has lagged policy
\begin{itemize}\setlength\itemsep{3pt}
\item Hidden by definition — actors actively invest resources to conceal it
\item Often associated with nonstate actors who capture state apparatus and media
\begin{itemize}\setlength\itemsep{3pt}
\item Thus, difficult to get accurate estimates
\item Reliance on corruption perception as measurement, but it can be very subjective
\end{itemize}
\item Theoretically difficult to model: Manifestation can vary in form across environments
\end{itemize}
\end{frame}

\begin{frame}
\frametitle{Corruption Perceptions Index: Little Progress Over Two Decades}
\begin{figure}
\centering
\includegraphics[width=0.75\textwidth]{fig3.png}
\caption*{\footnotesize Source: Transparency International. CPI scale: 0 (highly corrupt) to 100 (very clean).}
\end{figure}
\end{frame}


\begin{transitionframe}
\frametitle{Roadmap for today}
\tableofcontents
\end{transitionframe}








\section{Allocating authority across bureaucrats and politicians}


\subsection{Aghion, Tirole (1997): Formal and Real Authority in Organizations}
\begin{transitionframe}
\begin{center}
  \begin{huge}
    \textcolor{NTHU1}{%
Aghion, Tirole (1997):\\ Formal and Real Authority in Organizations
    }
  \end{huge}
\end{center}
\end{transitionframe}


\begin{frame}
\frametitle{Q: What decides the allocation of authorities within organizations?} 
What do we mean by authority
\begin{itemize}\setlength\itemsep{3pt}
\item Formal authority: Legal right to select actions (who owns the decision, final say)
\item Real authority: Effective control over decisions (who shapes the agenda)
\item  It is commonplace for these authorities to be held by different individuals
\begin{itemize}\setlength\itemsep{3pt}
\item Those who hold formal authority may delegate to better-informed agents
\item Shareholders/Presidents (formal) vs Board of directors/Executive branch (real)
\end{itemize}
\item Optimal decision-making structure when not everyone is equally informed?
\end{itemize}\medskip
This paper discusses optimal allocation and  separation of different types of authorities
\begin{itemize}\setlength\itemsep{3pt}
\item A principal delegates a task (project selection) to an agent
\item Both independently collect information about potential projects (at a cost)
\item The agent recommends a project; the principal approves or vetoes
\item Key question: When should the principal \textbf{retain} or \textbf{delegate} formal authority?
 \end{itemize}
\end{frame}






\begin{frame}
\frametitle{Model setup} 
Principal-agent can implement 0 or 1 among $n\geq 3$ projects
\begin{itemize}\setlength\itemsep{3pt}
\item Project $k$ has profit $B_k$ to principal and private benefit $b_k$ to agent
\item Principal's favorite project: $B$ to the principal and $\beta b$ to the agent $(\beta\in(0,1])$
\item Agent's favorite project: $b$ to the agent and $\alpha B$ to the principal  $(\alpha\in(0,1])$
\item Payoffs of the project initially unknown (perceived as identical)
\end{itemize}\medskip 
Utility and cost to principal and agent
\begin{itemize}\setlength\itemsep{3pt}
\item Principal pays wage $w$ to agent: Thus has utility $B_k-w$
\item Agent has income $w$ and private payoff $b_k$: Thus utility $u(w)+b_k$
\item Gathering information is costly (even if effort is exerted may not learn anything at all)
\begin{itemize}\setlength\itemsep{3pt}
\item Agent: Spend $g_A(e)$ to learn payoffs of all projects with probability $e$ 
\item Principal: Spend $g_P(E)$ to learn payoffs of all projects with probability $E$ 
\end{itemize}
\end{itemize}\medskip 
\end{frame}



\begin{frame}
\frametitle{Model setup} 
Timing of the decision structure
\begin{itemize}\setlength\itemsep{3pt}
\item Contract: Specifies wage and allocation of formal authority
\item Principal and agent simultaneously choose information acquisition efforts ($E$ and $e$)
\item Agent observes own information and recommends a project (or stays silent)
\item Principal observes own information and approves or vetoes the recommendation
\item If no project is approved, status quo payoff of 0 is realized
\end{itemize}\medskip
Utility (from project selection) depending on who has formal authority
\begin{itemize}\setlength\itemsep{3pt}
  \item Principal has formal authority (integration)\vspace{-10pt}
  \[
  u_P =  EB + (1-E)e\alpha B - g_P(E) \quad u_A = E\beta b + (1-E)eb-g_A(e)
  \]\vspace{-20pt}
  \item Agent has formal authority (delegation)\vspace{-10pt}
\[
  u_P^d =  e\alpha B + (1-e)EB - g_P(E) \quad u_A^d = e b + (1-e)E\beta b-g_A(e)
  \]\vspace{-20pt}
\end{itemize}
\end{frame}


\begin{frame}
\frametitle{Key trade-off: Principal's control vs. Agent's initiative} 
Efforts of principal and agents are \textit{strategic substitutes}: Principal's $E\uparrow\to$ Agent's $e\downarrow$
\begin{itemize}\setlength\itemsep{3pt}
  \item The first order conditions in the integration case is \vspace{-10pt}
  \[
  \text{P : } (1-e\alpha)B = g_P'(E) \quad \text{A : } (1-E)b=g_A'(e)
  \]\vspace{-20pt}
  \item FOC under delegation \vspace{-10pt}
  \[
  \text{P : } (1-e)B = g_P'(E) \quad \text{A : } (1-E\beta)b=g_A'(e)
  \]\vspace{-20pt}

  \item For instance, suppose that MC of searching increases in integrated case  ($g_p'\uparrow$)
  \begin{itemize}\setlength\itemsep{3pt}
  \item Probability of being informed about payoffs decreases $E\downarrow$ $\to$ Principal loses control
  \item But it can encourage agent to take initiative ($e\uparrow$)
\end{itemize}
\item You can argue similarly for delegation case (A loses initiative $\to$ P takes control)
\item Effort to collect information increases with the following factors..
  \begin{itemize}\setlength\itemsep{3pt}
  \item Principal: Lower congruence ($\alpha$), higher stakes ($B$), lower agent effort ($e$)
  \item Agent:  Lower congruence ($\beta$), higher private gains ($b$), lower principal interference ($E$)
\end{itemize}
\end{itemize}\medskip
\end{frame}




\begin{frame}
\frametitle{Optimal formal authority: When should principal delegate?}
Comparing equilibrium under integration $\{E, e\}$ and delegation $\{E^d, e^d\}$
\begin{itemize}\setlength\itemsep{3pt}
\item $e^d > e$: Agent exerts \textbf{greater initiative} under delegation
\item $E^d < E$: Principal exerts \textbf{less effort} — loses both formal and real authority
\item Agent more likely to find preferred project $\Rightarrow$ higher payoff $b_k$ $\Rightarrow$ accepts lower wage
\end{itemize}\medskip
Principal prefers to delegate over integration when initiative gains outweigh loss of control
\begin{itemize}\setlength\itemsep{3pt}
\item High congruence ($\alpha \uparrow$): Agent's preferred project less costly to principal
\item High private gains ($b \uparrow$): Agent more motivated to acquire information
\item High cost of principal's information acquisition ($g_P' \uparrow$): Retaining control is costly
\end{itemize}\medskip
Principal is fine with losing control if agents can find and implement good projects 
\begin{itemize}\setlength\itemsep{3pt}
  \item Especially when principal is overloaded, there are many superiors, and matter is urgent
\end{itemize}
\end{frame}

\begin{frame}
\frametitle{Span of control: Optimal \# of agents and degree of overload?}
Suppose that 
\begin{itemize}\setlength\itemsep{3pt}
  \item Principal retains formal authority 
  \item Hires $m$ identical agents, and monitoring them is costlier
\end{itemize}\medskip
Adding an extra agent has two opposing effects:
\begin{itemize}\setlength\itemsep{3pt}
\item Overload cost: Extra agent costs more in attention than revenue gained
\item Initiative effect: $e\uparrow$ $\Rightarrow$ agents find better projects $\Rightarrow$ principal's revenue increases
\end{itemize}\medskip
At the optimum, overload cost exactly offsets initiative effect
\begin{itemize}\setlength\itemsep{3pt}
\item Principal overhires beyond the point where marginal profit is zero 
\item The extra cost is made up for by gains from increased initiative
\item Busy principal cannot micromanage  $\to$ Overload as commitment not to stifle initiatives
\end{itemize}
\end{frame}

\begin{frame}
\frametitle{Key takeaways and discussions}
Other findings
\begin{itemize}\setlength\itemsep{3pt}
\item Participation view: delegate decisions that matter more to agent than principal 
\item Higher wages increase agent's real authority 
\begin{itemize}\setlength\itemsep{3pt}
  \item Stronger incentives $\to$ agent more likely to make recommendations
  \item High wage reduce principal's net profit $\Rightarrow$ Monitor less $\Rightarrow$ less likely to overrule agent
\end{itemize}
\item Integration may jeopardize communication 
\begin{itemize}\setlength\itemsep{3pt}
\item When congruence is low, agent more likely to withhold information
\end{itemize}
\end{itemize}\medskip
Open questions
\begin{itemize}\setlength\itemsep{3pt}
\item How does optimal authority allocation change with multilayered hierarchies?
\begin{itemize}\setlength\itemsep{3pt}
  \item Many middle managers or many levels of subnational government
  \end{itemize}
\item Authority or monetary incentives: what attracts and motivates agents?
\begin{itemize}\setlength\itemsep{3pt}
\item Agents whose decisions are not overruled may be more motivated to stay
\item But others may value monetary incentives over decision rights
\end{itemize}
\end{itemize}
\end{frame}




\subsection{Bandiera, Best, Khan, Prat (2021): The Allocation of Authority in Organizations: A Field Experiment with Bureaucrats}
\begin{transitionframe}
\begin{center}
  \begin{huge}
    \textcolor{NTHU1}{%
Bandiera, Best, Khan, Prat (2021):\\ The Allocation of Authority in Organizations: A Field Experiment with Bureaucrats
    }
  \end{huge}
\end{center}
\end{transitionframe}




\begin{frame}
\frametitle{Q: How does allocation of authority to bureaucrats affect outcomes?} 
Organization bring together people with different interests and skills
\begin{itemize}\setlength\itemsep{3pt}
\item Working towards a common goal challenging: Principal-agent problem and monitoring
\item How to allocate decision rights at different layers of hierarchy?
\item How to monitor and incentivize agents' behavior?
\item Many work on effects of performance pay, few look at decision-making structure
\end{itemize}\medskip
This paper designs an experiment in Pakistan procurement sector that combines the two
\begin{itemize}\setlength\itemsep{3pt}
\item Strict rules and monitoring vs. simplification and more autonomy?
\item Procurement officers (PO, make purchases) and Monitors (approve purchase) 
\begin{itemize}\setlength\itemsep{3pt}
\item Monitors: Approve purchases quickly vs Delay purchase until end of fiscal year
\end{itemize}
\item Experiment: Control + Autonomy + Performance pay + Both 
\item Different policy instrument works in different structures
\end{itemize}
\end{frame}

\begin{frame}
\frametitle{Procurement structure in Pakistan and theoretical framework} 
Procurement structure
\begin{itemize}\setlength\itemsep{3pt}
\item Experiment took place in 26 / 36 districts in Punjab (80\% of population)
\item Each office has 1 PO with allocated budgets
\item POs submit purchase for pre-audit approval to a monitor in each district
\begin{itemize}\setlength\itemsep{3pt}
\item This monitor can ask for extra paperwork (withhold) or approve the purchase
\end{itemize}
\end{itemize}\medskip
Conceptual framework: Expected effects of autonomy vs incentives by type of monitors
\begin{itemize}\setlength\itemsep{3pt}
\item Aligned PO and monitors behave in the interest of citizens (low price, little delays)
\item Misaligned PO raises price (bribes, little effort to find low-price)
\item Misaligned monitor delays purchase (exert power on PO)
\item Autonomy: Price$\downarrow$ if PO more aligned then monitor (opposite if PO more misaligned)
\item Incentives: Motivate PO, but dampened with misaligned monitors (higher costs)
\end{itemize}
\end{frame}

\begin{frame}
\frametitle{Experimental design} 
Measuring bureaucratic performance with value for money (with online platform)
\begin{itemize}\setlength\itemsep{3pt}
\item Price conditional on quantity, nature of purchased goods, and net of transaction costs
\begin{itemize}\setlength\itemsep{3pt}
\item Items: Consumables such as papers, toners, pen, banners, light bulbs etc
\item Price is adjusted for delivery speed, cost of delivery
\end{itemize}
\item Others: Expected price, good-bad purchase categorization(based on attributes), ML
\end{itemize}\medskip
Experimental design
\begin{itemize}\setlength\itemsep{3pt}
\item Autonomy: Remove the role of monitors in purchasing decisions
\item Incentive: Range of bonus for improving value for money ($\times2$  - $\times1/2$ of monthly wage)
\item Another group that receives both
\item Stratified by district-department level\vspace{-10pt}
 \[
 y_{igto} = \alpha + \sum_{k=1}^3\eta_k \text{Treatment}_o^k + \rho_g q_{igto}+\delta_s +\gamma_g + e_{igto}
 \]
\end{itemize}
\end{frame}

\begin{frame}
\frametitle{Overall, autonomy $\to$ lower prices paid but not incentives} 
\begin{figure}
  \includegraphics[keepaspectratio, width=\textwidth]{bbkp2021_t1.png}
\end{figure}
\end{frame}

\begin{frame}
\frametitle{Effects differ by type of monitors} 
\begin{figure}
  \includegraphics[keepaspectratio, width=0.7\textwidth]{bbkp2021_f1.png}
\end{figure}
\begin{itemize}\setlength\itemsep{3pt}
\item Better monitors (low delays, AG June share $\downarrow$): Incentive lower price (not autonomy)
\item Worse monitors (more delays, AG June share $\uparrow$): Autonomy lower price (not incentives)
\end{itemize}
\end{frame}

\begin{frame}
\frametitle{Treatment $\to$ large savings, but heterogenous by monitor types} 
\begin{figure}
  \includegraphics[keepaspectratio, width=0.7\textwidth]{bbkp2021_f2.png}
\end{figure}
\end{frame}

\begin{frame}
\frametitle{Takeaways and remaining questions} 
Incentives and autonomy work in different situations
\begin{itemize}\setlength\itemsep{3pt}
\item Incentives to low-level agents mean little if their supervisors are misaligned
\item In such case, giving low-level agents more autonomy improves performance
\item If supervisors are aligned, incentives improve effort of angents
\item On the other hand, autonomy does not change much from status quo in this case
\end{itemize}\medskip
Remaining questions on the organizational Economics
\begin{itemize}\setlength\itemsep{3pt}
\item Can stringent rules, not autonomy, be an optimal choice?
\begin{itemize}\setlength\itemsep{3pt}
\item Could still be worthwhile to prevent misaligned agents imposing high costs
\end{itemize}
\item How does autonomy affect \textit{who becomes a bureaucrat}?
\begin{itemize}\setlength\itemsep{3pt}
\item How to bring in more aligned agents (make efficient decisions for public) than misaligned agents (extract personal benefits with autonomy)?
\end{itemize}
\end{itemize}\medskip
\end{frame}


\section{Analysis of Corruption}
\subsection{Banerjee, Mullainathan, Hanna (2012): Corruption}
\begin{transitionframe}
\begin{center}
  \begin{huge}
    \textcolor{NTHU1}{%
Banerjee, Mullainathan, Hanna (2012):\\ Corruption
    }
  \end{huge}
\end{center}
\end{transitionframe}


\begin{frame}
\frametitle{Q: What are the causes and consequences of corruption?}
Despite rampant corruption, research has lagged reality
\begin{itemize}\setlength\itemsep{3pt}
\item Corruption is empirically difficult to measure
\item Theoretical analysis focused exclusively on moral hazard in organizations
\begin{itemize}\setlength\itemsep{3pt}
\item Emphasis on corruption as a criminal action requiring monitoring and strict penalties 
  \item In practice, corruption can manifest in many different forms across different environment
\end{itemize}
\item Need to understand corruption from a different standpoint
\end{itemize}\medskip
This paper provides a novel framework for analyzing corruption in public organizations
\begin{itemize}\setlength\itemsep{3pt}
\item Corruption as result of tasks, private incentives, and organizational environment
\item Highlight how discretion (vs. rules) affects corruption in different environments
\item Also study other distortions that ccur with corruption 
\begin{itemize}\setlength\itemsep{3pt}
\item Red-tape, inefficiencies in service provision
\end{itemize}
\end{itemize}
\end{frame}



\begin{frame}
\frametitle{What is corruption} 
\textbf{Rule-breaking} by a bureaucrat (or elected official) for \textbf{private gain}
\begin{itemize}\setlength\itemsep{3pt}
\item Overt monetary bribe to bend a rule and provide service to unintended recipient
\item Nepotism as opposed to competitive hiring / selection
\item Not showing up to work (`steal' time)

\item These are defined by rules specific to a context
\end{itemize}\medskip
How this definition differs from other approaches
\begin{itemize}\setlength\itemsep{3pt}
\item Emphasis on breaking the rules (previous definitions emphasize personal gain)
\begin{itemize}\setlength\itemsep{3pt}
\item A boss \textit{with final say} hiring his nephew: Not corruption if no other rules broken
\item Previous definitions would count this as corruption (personal gain)
\end{itemize}
\item Advantages of this definition:
\begin{itemize}\setlength\itemsep{3pt}
\item Avoids subjective ethical judgements — rules define what is corrupt
\item Encompasses non-monetary gains (nepotism, time theft)
\item Rule-breaking is easier to measure than bribes
\end{itemize}
\end{itemize}
\end{frame}



\begin{frame}
\frametitle{Model setup}
A bureaucrat assigns a limited number of slots to a population of size $N$
\begin{itemize}\setlength\itemsep{3pt}
\item Individuals differ in social value, private value, and ability to pay
\item Individuals announce their types (private value is hidden)
\end{itemize}\medskip
Bureaucrat announces mechanism specifying testing, prices, and assignment probabilities
\begin{itemize}\setlength\itemsep{3pt}
\item Testing: screen individual's types at a cost to both bureaucrat and individual
\item Assignment: probability of receiving a slot given test outcome
\item Prices: what individuals pay for a slot
\end{itemize}\medskip
Government sets rules that constrain the mechanism (and define what is corrupt)
\begin{itemize}\setlength\itemsep{3pt}
\item If rules don't pin down a unique mechanism $\Rightarrow$ bureaucrat has \textbf{discretion}
\item Various dimensions of malfeasance beyond corruption
\begin{itemize}\setlength\itemsep{3pt}
  \item Overcharging price (bribe), excessive testing (red-tape), no testing (shirking)
  \item Misallocation (allocative inefficiency): slots assigned to wrong types or left unassigned
\end{itemize}
\end{itemize}
\end{frame}


\begin{frame}
\frametitle{The task approach to corruption}
Corruption arises from the nature of the task itself
\begin{itemize}\setlength\itemsep{3pt}
\item Same bureaucrat, same incentives — different tasks generate different corruption
\item Can we eliminate corruption by changing the rules without affecting anything else?
\end{itemize}\medskip
Key tradeoff between setting strict rules and allowing discretion
\begin{itemize}\setlength\itemsep{3pt}
\item Strict rules limit corruption but generates red tape and misallocation
\begin{itemize}\setlength\itemsep{3pt}
\item Could ensure that public services are allocated based on needs (welfare)
\item Strict rules bind honest bureaucrats but corrupt ones find ways around them
\end{itemize}
\item Full discretion reduces red tape but opens door to corruption
\begin{itemize}\setlength\itemsep{3pt}
\item Italian procurement: Fixed budget with full discretion reduced waste and corruption
\item Hospital bed allocation: May lead to allocation to who can pay, not need-based
\end{itemize}
\end{itemize}\medskip
Optimal response: redesign the task (rules \& discretion), not just increase punishment
\begin{itemize}\setlength\itemsep{3pt}
\item Goal is to improve public service provision (punishment is one of the tools)
\end{itemize}
\end{frame}

\begin{frame}
\frametitle{Manifestations of corruption}
The form corruption takes depends on rules and gap between social and private values
\begin{itemize}\setlength\itemsep{3pt}
\item Bribes: Bureaucrat has pricing discretion $\Rightarrow$ overcharges those who can pay
\begin{itemize}\setlength\itemsep{3pt}
\item Arises when private value exceeds social value — bureaucrat exploits the gap
\end{itemize}
\item Red tape: Bureaucrat cannot price freely $\Rightarrow$ imposes excessive testing costs
\item Misallocation: Bureaucrat has assignment discretion $\Rightarrow$ slots go to those who can pay
\begin{itemize}\setlength\itemsep{3pt}
\item Arises when social value is high but ability to pay is low (e.g. poor patient, small farmer)
\end{itemize}
\end{itemize}\medskip
These distortions arise simultaneously and reducing one dimension may worsen others
\begin{itemize}\setlength\itemsep{3pt}
\item Restrict pricing $\Rightarrow$ Bribes fall but red tape rises
\item Restrict testing $\Rightarrow$ Red tape falls but misallocation rises
\item Restrict assignment $\Rightarrow$ Misallocation falls but bribes rise (value in bending rules)
\item Optimal policy: redesign the task so that corruption causes the least harm
\end{itemize}
\end{frame}



\begin{frame}
\frametitle{Measuring corruption: Empirical strategies used in the literature}

Perception-based measures (old approach): surveys of participants
\begin{itemize}\setlength\itemsep{3pt}
\item Problems: subjective, culturally biased, hard to compare across settings
\end{itemize}\medskip
Objective measures (new approach): direct measurement or administrative data
\begin{itemize}\setlength\itemsep{3pt}
\item Price comparisons (procurement): Price discrepancies for identical goods
\item Audit studies: Send trained auditors to measure bribes and red tape directly
\item Absenteeism: Photograph or timestamp bureaucrat presence — captures time theft
\item Leakage studies: Track whether public resources actually reach intended recipients
\end{itemize}\medskip
Key challenge: Measurement itself changes behavior (Hawthorne effect)
\begin{itemize}\setlength\itemsep{3pt}
\item Actors reduce illicit behavior during measurement or find new ways to obscure it
\end{itemize}
\end{frame}


\begin{frame}
\frametitle{Policy Implications and Open Questions}
What works against corruption?
\begin{itemize}\setlength\itemsep{3pt}
\item Monitoring / punishment: Effective but limited — bureaucrats adapt and hide behavior
\item Task redesign: Change rules and discretion to minimize the most costly distortion
\item Wages: Higher wages may reduce corruption by raising cost of job loss
\item Competition: Multiple providers reduces bureaucrat's monopoly over slots
\item Nonmonetary benefits: Attract intrinsically motivated, prosocial bureaucrats
\end{itemize}\medskip
Open questions
\begin{itemize}\setlength\itemsep{3pt}
\item How to recruit intrinsically motivated individuals to public sector?
\begin{itemize}\setlength\itemsep{3pt}
\item Public sector competes with private sector
\item Higher wages \textit{could} work, but may also attract undesirable individuals
\end{itemize}
\item How does allowing discretion deter corruption in practice?
\begin{itemize}\setlength\itemsep{3pt}
\item Could increase initiative and overrule corruptible middle managers
\item But also increase scope for corruption due to less interventions
\end{itemize}
\end{itemize}
\end{frame}




\subsection{Dal Bó, Dal Bó, and Di Tella (2006): Plata o Plomo? Bribe and Punishment in a Theory of Political Influence}
\begin{transitionframe}
\begin{center}
  \begin{huge}
    \textcolor{NTHU1}{%
Dal Bó, Dal Bó, and Di Tella (2006):\\ Plata o Plomo? Bribe and Punishment in a Theory of Political Influence
    }
  \end{huge}
\end{center}
\end{transitionframe}


\begin{frame}
\frametitle{Q: How do state capture affect quality of public officials} 
In weak states, politicians are subject to punishments by private entities and bribes
\begin{itemize}\setlength\itemsep{3pt}
  \item Many private entities seek to influence politicians through bribes and coercion
\item Punishments and bribes complement: Threat of punishment makes bribing easier
\item Yet, there are things where models with only corruption cannot explain
\begin{itemize}
\item Why give officials more discretion and immunity?
\end{itemize}
\end{itemize}\medskip
This paper updates framework of political influence to include bribes and coercion
\begin{itemize}\setlength\itemsep{3pt}
\item Investigates how possibility of coercion affects quality of public workers 
\item Explains why political immunities exist in contexts where private coercion is prevalent
\item Explains how discretion actually makes public work less attractive
\begin{itemize}\setlength\itemsep{3pt}
\item In model with bribes only, immunities are rarely justified (weaken monitoring mechanisms)
\item Also in the old model, discretion increases corruption and appeal of public work
\end{itemize}
\end{itemize}
\end{frame}

\begin{frame}
\frametitle{Model setup: Decision to join public work}
First stage: Citizens decide to enter public office or not 
\begin{itemize}\setlength\itemsep{3pt}

\item Citizens endowed with ability $a\in[0,\infty)$ (a probability distribution)
\item Alternative is to choose private sector (which gives wage $a$, equivalent to their ability)
\item Return from public work: public sector wage ($w$) + behavior of pressure group 
\item Citizens with $a\leq w$ apply (and $w$ types selected, assuming recruiters observe types)
\end{itemize} \medskip
Second stage: Pressure group has an opportunity to influence officials
\begin{itemize}\setlength\itemsep{3pt}
\item Public workers provide public good, whose quality positively correlates with ability
\item Have discretion to allocate ($\pi$) towards pressure group, with cost $c$ if caught
\item With probability $\gamma$, officials decide to accept bribes or face punishments
\begin{itemize}
  \item $1-\gamma$: Only punishment (no possibility to accept bribes) - when bribes are infeasible
\end{itemize}
\item Pressure groups can either bribe ($b$) or punish ($p$) at costs $\beta\Phi(b)$ and $\rho\Psi(p)$
\end{itemize}
\end{frame}

\begin{frame}
\frametitle{Analysis of payoffs} 
Decision on accepting (and handing out) bribes
\begin{itemize}\setlength\itemsep{3pt}
\item Officials facing pressure groups accept bribes (share $\gamma$) if \vspace{-5pt}
\[
w+b-c \geq w-p
\]
\vspace{-20pt}
\item Pressure groups sets $b$ and $p$ to maximize the following profits\vspace{-5pt}
\[
\Pi(b,p)=\gamma(\pi-\beta\Phi(b))-(1-\gamma)\rho\Psi(p) \quad \text{s.t.  }\quad b\geq c-p
\]
\vspace{-20pt}
\item $\tilde{\pi}$: Minimal discretion required for pressure groups to influence officials
\begin{itemize}\setlength\itemsep{3pt}
\item $\pi$ that makes $\Pi(b^*,p^*)=0$ ($\pi$ that makes profits when groups optimize $0$)
\item Inversely related to state capture: $\tilde{\pi}\Uparrow$ $\to$ Cost of bribing an official $\Uparrow$ $\to$ State capture  $\Downarrow$
\end{itemize}
\end{itemize}\medskip
If only bribes are available ($p=0$ for all cases)
\begin{itemize}\setlength\itemsep{3pt}
\item Solve $\gamma(\pi-\beta\Psi(b))$ subject to $b\geq c$ $\to$Bribe just enough to cover $c$, ($b=c$)
\item Official payoffs/skill is $w$, and $\tilde{\pi}^0=\beta\Phi(c)$
\item Drop in bribe costs ($\beta\downarrow$) $\to$ Higher state capture ($\tilde{\pi}\Downarrow$), no effect on quality of officials
\end{itemize}
\end{frame}

\begin{frame}
\frametitle{What happens when both bribes and coercion are possible?} 
Problem when both $b$ and $p$ are active
\begin{itemize}\setlength\itemsep{3pt}
\item Pay bribes just enough to cover $b=c-p$ (Bribes become cheaper)
\item The shortfall is made up with nonzero punishment $p^*$, such that $b+p^*=c$
\item $\tilde{\pi}^p$ satisfies $\frac{1-\gamma}{\gamma}\rho\Psi(p^*)+\beta\Phi(c-p^*)$ $\to$ Any discretion $\pi$ larger than this, then influence 
\item Resulting equilibrium payoff of official if pressure groups are active
\begin{itemize}\setlength\itemsep{3pt}
\item If official refuses bribe, payoff is $w-p^*$ (lower due to nonzero punishments)
\item If bribes are accepted, payoff is $w+b-c = w+c-p^* -c = w-p^*$
\end{itemize}
\item When pressure groups are inactive: payoff is $w>w-p^*$
\item Thus, the quality of bureaucrats fall in this scenario
\item Furthermore, $\tilde{\pi}^p<\tilde{\pi}^0$ $\to$ When threats are usable, more state capture!  
\end{itemize}
\end{frame}

\begin{frame}
\frametitle{Negative effects amplified with cheaper ways to influence officials} 
Exogenous changes in cost of influence (Suppose that $\beta$ and $\rho$ drops)
\begin{itemize}\setlength\itemsep{3pt}
\item Recall that $\tilde{\pi}^p=\frac{1-\gamma}{\gamma}\rho\Psi(p^*)+\beta\Phi(c-p^*)$ and $\Pi(b,p)=\gamma(\pi-\beta\Phi(b))-(1-\gamma)\rho\Psi(p)$
\item Fall in any one of $\beta$  and $\rho$ raises profits and decreases $\tilde{\pi}^p$
\begin{itemize}\setlength\itemsep{3pt}
\item State capture will be more rampant 
\item Existing pressure groups ramp up activity + new pressure groups enter
\end{itemize}
\item If primarily driven by $\rho\Downarrow$, then payoff of officials (and quality) definitely falls
 \begin{itemize}\setlength\itemsep{3pt}
\item Usage of coercion increases in any case
\end{itemize}
\item If primarily driven by $\beta\Downarrow$, change to pay off ambiguous (but definitely not increase)
 \begin{itemize}\setlength\itemsep{3pt}
\item Existing groups shift away from coercion $\to$ Minimal effect on wages
\item New entrants that use both coercion and bribes $\to$ lower payoffs

\end{itemize}
\end{itemize}
\end{frame}

\begin{frame}
\frametitle{Implication on discretions}
How do discretion affect quality of politicians 
\begin{itemize}\setlength\itemsep{3pt}
\item If discretion $\pi<\tilde{\pi}^p$, less incentives for pressure groups to influence officials
\item Higher discretion could lead to more waste (more corrupt decision)
\item As more groups with coercion decrease payoffs, quality of officials worsened
\end{itemize}\medskip
Immunity of officials (in principle)
\begin{itemize}\setlength\itemsep{3pt}
\item Can insulate honest officials from outside pressures, mitigate false accusations
\item But can mask corrupt officials
\item In setting with rampant coercion from pressure groups
\begin{itemize}\setlength\itemsep{3pt}
\item There are greater effect of protection to officials with immunity
\item This can overweigh bad effects of masking corrupt officials
\end{itemize}
\item When legal system is effective, masking effects outweigh pressure effects
\end{itemize}
\end{frame}

\begin{frame}
\frametitle{Takeaways and discussions} 
Threat of coercion imposes negative effects on governance
\begin{itemize}\setlength\itemsep{3pt}
\item Reduces payoff for officials $\to$ lower quality of officials
\item More influence from private pressure groups $\to$ Increased state capture
\item Augments effects of bribes: Cheaper to bribe officials
\end{itemize}\medskip
Remaining questions
\begin{itemize}\setlength\itemsep{3pt}
\item Ultimately, need to improve legal capacity 
\begin{itemize}\setlength\itemsep{3pt}
\item To prevent pressure groups and detect corruption
\item But transitioning to better legal system is tough
\end{itemize}
\item Incentives of bureaucrats: How to attract incorruptible moral officials?
\begin{itemize}\setlength\itemsep{3pt}
\item Under what conditions can governments bring them in (what sort of incentives?)
\item Topic of Dal Bó et al (2013)
\end{itemize}
\end{itemize}\medskip
\end{frame}



\subsection{Xu (2018): The Costs of Patronage: Evidence from the British Empire}
\begin{transitionframe}
\begin{center}
  \begin{huge}
    \textcolor{NTHU1}{%
Xu (2018):\\ The Costs of Patronage: Evidence from the British Empire
    }
  \end{huge}
\end{center}
\end{transitionframe}

\begin{frame}
\frametitle{Q: How does patronage affect allocation of positions \& performance?} 
How personnel are promoted and incentivized determines organizational performance
\begin{itemize}\setlength\itemsep{3pt}
\item Past: patronage widely used to appoint public sector positions
\begin{itemize}
\item Discretionary appointments based on network and ties to the ruler
\end{itemize}
\item In theory, discretion can align incentives between principal and angents
\item However, it can worsen performance if favoritism disincentivizes capable subordinates
\end{itemize}\medskip
This paper uses historical data to identify how patronage affects performance
\begin{itemize}\setlength\itemsep{3pt}
\item Colonial governors of British empire: 1854-1966 across 70 colonies
\item Connection determined by genealogy and biological data
\item Within-governor variation: Change in secretary of state $\to$ Change in connections
\item Across-time variation: Discretion (1854-1930) vs. Independent appointment (1930 -)
\item Test who gets assigned to richer colonies and how tax returns are affected
\end{itemize}

\end{frame}

\begin{frame}
\frametitle{Organization of British colonies}
Assignment of colonial governors 
\begin{itemize}\setlength\itemsep{3pt}
\item Handled by Colonial Office under Secretary of State for the Colonies in London
\item Secretary of State in charge for around 3 years
\item 1930 Reform replace Secretary of State's discretion with open recruitment
\end{itemize}\medskip
Role of the governors
\begin{itemize}\setlength\itemsep{3pt}
\item Fixed 6-year term
\item Responsible for public finance, legislation, and personnel appointments
\item Discretion over executive and finances: Widely different policies across colonies
\item Direct monitoring impossible due to distance
\end{itemize}\medskip
\end{frame}

\begin{frame}
\frametitle{Data and empirical strategy}
Digitalized (by Guo Xu himself!) archival data 
\begin{itemize}\setlength\itemsep{3pt}
\item Sources: Colonial Blue Books, Colonial Office List, Who's Who, The Peerage
\item Postings (and their wages), background of governors, genealogy 
\begin{itemize}
\item Connection determined by ancestry, membership at aristocratic clubs, and schools
\end{itemize}
\item Colonial statistics (revenue, expenditure, demographics etc)
\end{itemize}\medskip
Empirical strategy leverages differences in connectedness across different periods\vspace{-10pt}
\[
y_{it} = \beta c_{it}+\theta_i + \gamma X_{it}+\tau_t + e_{it}
\]\vspace{-15pt}
\[
y_{it} = \beta_0 c_{it} + \beta_1 c_{it}\times I[t\geq1930]+\theta_i + \gamma X_{it}+\tau_t + e_{it}
\]\vspace{-10pt}
\begin{itemize}\setlength\itemsep{3pt}
\item Connection $c_{it}$ used with various criterion
\item $c_{it}$ is predetermined and driven by secretary of state changes
\begin{itemize}
  \item No reverse causality: Unlikely to attend Oxbridge to be connected to current governor
  \item Secretary of state turnover generates common shocks to all serving governors
\end{itemize}
\end{itemize}
\end{frame}

\begin{frame}
\frametitle{Connected governor salaries are higher! } 
    \begin{figure}
    \includegraphics[keepaspectratio, width=0.8\textwidth]{x2018_t1.png}
    \end{figure}
\end{frame}

\begin{frame}
\frametitle{Salary gap washed away post-1930} 
    \begin{figure}
    \includegraphics[keepaspectratio, width=0.7\textwidth]{x2018_f1.png}
    \end{figure}
\end{frame}

\begin{frame}
\frametitle{Patronage worsens performance, which is offset following reforms} 
\begin{columns}
  \begin{column}{0.65\textwidth}
    \begin{figure}
    \includegraphics[keepaspectratio, width=\textwidth]{x2018_t2_1.png}
    \end{figure}
  \end{column}
  \begin{column}{0.3\textwidth}
    \begin{figure}
       \includegraphics[keepaspectratio, width=\textwidth]{x2018_t2_2.png}
       \end{figure}
  \end{column}
\end{columns}
\end{frame}

\begin{frame}
\frametitle{Takeaways and discussions} 
What the findings say about organizational performance
\begin{itemize}\setlength\itemsep{3pt}
\item Connection $\to$ Positions that yield higher wages
\item Lower revenue collection: Due to higher tax exemptions, high reliance on customs
\item Reforms washed salary gaps and shortfalls in performance
\end{itemize}\medskip
Discussions
\begin{itemize}\setlength\itemsep{3pt}
\item Favoritism can lead to worse organizational performance, even in public institutions
\item But how does it affect recruitment into the pool of governor candidates?
\begin{itemize}\setlength\itemsep{3pt}
\item If people know connections matter, talented individuals with no connection discouraged
\item Leaving pool of candidates with less talented individuals
\item Ways to look for `outsiders' with talent?
\end{itemize}
\end{itemize}
\end{frame}


\begin{frame}
\frametitle{Remaining questions} 
How to attract talented individuals to the public sector?
\begin{itemize}\setlength\itemsep{3pt}
\item In a lot of places, they have lower pay
\item But public sector offers different perks: Stability, social respect, instrinsic motivation
\item How can these individuals be recruited?
\begin{itemize}\setlength\itemsep{3pt}
\item Does financial incentives work?
\item Are these individuals different from the general public?
\end{itemize}
\end{itemize}\medskip
How to design a healthy public organizations?
\begin{itemize}\setlength\itemsep{3pt}
\item Healthy balance of incentives and authority that prevents misallocation and corruption
\item Monitoring and incentives are still often used
\item But are there other aspects of public organizations that can be used to improve performance?
\begin{itemize}\setlength\itemsep{3pt}
\item Career progression: Initial and subsequent posting to different positions
\item Middle managers: Can they bridge the gap between agents and the principal (the people)
\end{itemize}
\end{itemize}
\end{frame}



%%%%%%%%%%%
\end{document}
